\cotefile{161-2}
\scanpage{96}
La catégorie (ou \unsure{un} \uline{lieu} quelc.) universelle
pour ceci est la catégorie des faisceaux sur
$R'$ pour la top. de Zariski (= foncteurs
contravariants $R\longrightarrow\mathrm{Ens}$ qui transforment
sommes en produits). Il se trouve que
c'est un topos \deleted{$\unsure{E|\mathcal E|\circledast}$} Le corps universel est
le faisceau structural évident sur ce topos....

\centerline{\rule{5em}{0.4pt}}

$\overline{\mathrm M}$ question avec \uline{anneau intègre}. Soit
$O'\subset O$ le sous-objet ``qui opère régulièrement sur $O$'',
$(O'\supset O^*)$ est il définissable par $\varprojlim$ et $\varinjlim$ ? Commute-t'il aux images inverses \unsure{par un morphisme} de topos ?

\deleted{$O'$ définissable par $\varprojlim$ \unsure{finies} ? Si $\unsure{E}=\unsure{\mathrm{Sch.aff}|\mathrm{Sch.ff}}^{\unsure{\sim}}$, cela
signifie que les foncteurs \unsure{fonctionnant} sur Sch aff
Cov. sur \unsure{Ann} qui : \unsure{tt}}

Dans le cas \unsure{de|où|des} l'anneau universel d'un
topos, $O\in\widehat R$, $O'(\deleted{\unsure{A|X}})$ est le sous-\deleted{\unsure{anneau}}\inserted{ens de}
de $A=O(X)$ formé des $f\in A$ tels que $f$ soit
régulier, et \uline{le reste} par tt chgt d'anneau
de type fini $A\longrightarrow A'$ -- mais cela \deleted{signifie}
implique que $f\notin m$ pour $m\in\mathrm{Max}\,A$, donc
$f$ inversible. Donc \unsure{on a} $O'=O^*$ ! Or
un foncteur fibre sur $\widehat R$ est un anneau
(ensembliste) et il n'est pas vrai qu'un tel
anneau soit intègre \unsure{ssi} c'est un corps, donc
on n'a pas $O'_p=O_p^{*}$ i.e $A'=A^*$ !

% archive sheet number: 95
\scanpage{98}
\mspage{\textcircled{1}}
\fbox{1} Soit $(1.1)\ \lambda=(\underset{\leftarrow}{\lambda},\underset{\rightarrow}{\lambda})$ un couple, où $\underset{\leftarrow}{\lambda}$ et $\underset{\rightarrow}{\lambda}$ sont des
ens. de diagrammes de l'Univers fixé $\mathcal U$.
On veut préciser la notion de ``structure \unsure{algébrique|algébriques} définie
\deleted{par les} à l'aide des $\varprojlim$ de type $\underset{\leftarrow}{\lambda}$, des $\varinjlim$ de type $\underset{\rightarrow}{\lambda}$''.
(ou $\lambda$-type) \deleted{\illegible{1}})

{\scriptsize
\[
\begin{aligned}
(\mathrm{Cat}_{\lambda})&=\text{2-catégorie }
\left\{\begin{aligned}
 &\underline{\text{objets}}: \inserted{\mathcal U\text{-}}\text{catégories ayant }
 \left\{\begin{aligned}&\varinjlim\ \text{type }\underset{\rightarrow}{\lambda}\\[-1pt]&\varprojlim\ \text{type }\underset{\leftarrow}{\lambda}\end{aligned}\right\}=\lambda\text{-catégories}\\
 &\text{\deleted{des catégories}}\\
 &\underline{\text{morphismes}}:\text{ foncteurs qui commutent aux}\\
 &\qquad\text{types de }\varprojlim\text{ et }\varinjlim\qquad\Big\} \lambda\text{-}\unsure{\text{fonct.}|\text{foncteurs}}\\
 &\underline{\text{morphismes}}\text{ de morphismes : comme dans }(\mathrm{Cat})
\end{aligned}\right.
\end{aligned}\tag*{(1.2)}
\]
}
\[
[\text{donc }\mathrm{Cat}_{\lambda}\hookrightarrow(\mathrm{Cat})\text{ est }\unsure{1|2}\text{-fidèle}]
\]

Si $T$ est une \deleted{espèce de structure} \deleted{$\lambda$-structure} $\lambda$-type, on aura
pour tt $C\in\mathrm{Ob}\,\mathrm{Cat}_{\lambda}$ \deleted{\unsure{la|les}} catégorie $T(C)$ des
structures de type $T$ dans $C$, et on trouve ainsi une
catégorie cofibrée sur $\mathrm{Cat}_{\lambda}$, soit $T_{\lambda}$
\[
\mathbf T_{\lambda}\text{ cat cofibrée sur }\mathrm{Cat}_{\lambda}\tag*{(1.3)}
\]
Si $T'$ est un autre \deleted{structure} $\lambda$-type, on pose
\[
\underline{\mathrm{Hom}}_{\lambda}(T,T')=\underline{\mathrm{Hom}}_{\mathrm{fib\ sur}\,(\mathrm{Cat}_{\lambda})}(T_{\lambda},T'_{\lambda})\tag*{(1.4)}
\]
et de cette façon les $\lambda$-types forment une (2-cat)
\[
(\lambda\text{-types})\overset{\text{2-fidèle}}{\lhook\joinrel\longrightarrow}\underline{2\text{-}\mathrm{Catfib}}\,(\mathrm{Cat}_{\lambda})\tag*{(1,5)}
\]

\fbox{2} Si on a donné $T$ avec ens d'indices $I$, tel que
``$T$ structure sur $I$ objets de base'', on introduit la
théorie $\tau_{\deleted{\unsure{I}}}^{I}$ (structure vide sur $I$ ens de base)
donnant lieu à la catégorie cofibrée sur $(\mathrm{Cat}_{\lambda})$, de fibres les
\[
\tau^I(C)=C^I\tag*{(2.1)}
\]
et on trouve \deleted{\unsure{des trois|des}} \inserted{pour tt $C$ un} foncteur canonique ``syst des objets
de base''
\[
T_{\lambda}(C)\xrightarrow{b_{\lambda}^{C,T}}C^{I}\qquad
\bigl[(\text{foncteur }\underline{\text{0-fidèle}}\,\inserted{\text{et }\unsure{\text{tr transportable}|\text{transportable}|\text{transposition}}})\bigr]\tag*{(2,2)}
\]
définissant pour $C$ variable
\[
b_{\lambda}^{\unsure{T|I}}:T_{\lambda}\longrightarrow\tau^{I}\qquad(\text{0-fidèle})\tag*{(2.3)}
\]
% archive sheet number: 97
\scanpage{99}
\mspage{\textcircled{2}}
On définit, si $T'$ est de $=$ relatif à $I$
\[
\begin{aligned}
\underline{\mathrm{Hom}}_\lambda(I;T,T')={}&\text{catégorie des }\inserted{\text{couples }(u,\alpha),\ \text{avec }u}\\
 &\text{morphismes } u:T_\lambda\longrightarrow T'_\lambda\ \text{\deleted{avec} isom de}\\
 &\text{commutation }\alpha
\end{aligned}\tag*{(2.4)}
\]
\diagram{\begin{tikzcd}[ampersand replacement=\&,column sep=small]
 T_\lambda \arrow[r,"u"]\arrow[d,"b_\lambda^T"']\& T'_\lambda\arrow[dl,"b_\lambda^{T'}"]\\
 \tau^I \& \deleted{\unsure{\mathrm{id}}}
\end{tikzcd}}

On a une \deleted{foncteur} 2-cat.
\[
(\lambda\text{-}I\text{-types})\lhook\joinrel\longrightarrow
\unsure{\text{(}}\underline{\mathrm{Cat}\ 2\,\mathrm{cofib}}\,(\mathrm{Cat}_{\lambda}))/\tau^{I}\tag*{(2.5)}
\]
et un 2-foncteur
\[
(\lambda\text{-}I\text{-types})\longrightarrow(\lambda\text{-types})\tag*{(2.6)}
\]
qui est $\unsure{0|1}$-fidèle.

\fbox{3} Si $I$ est réduit à un élément, $\tau^I$ est noté \deleted{$\unsure{I}$} $\tau$, donc
\[
\tau(C)=C\tag*{(3.1)}
\]
\deleted{\illegible{1} catégorie} On pose
\[
\underline{\mathrm{Hom}}_\lambda(T,\tau)\simeq R_T^\lambda\quad\deleted{\unsure{\mathrm{Rep}|(\text{ou }R_T)}}\tag*{(3,2)}
\]
et l'appelle Cat. des $\lambda$-\uline{constructions pour} $T$. On \deleted{\unsure{trouve}}\inserted{ne}
\deleted{\unsure{peut}} sait pas à priori si $R_T^\lambda$ est une $\mathcal U$-catégorie. Mais
c'est une \deleted{\unsure{cat}} $\lambda$-catégorie, et les $\lambda$-limites se calculent
``argument par argument'', donc pour toute $C\in\mathrm{Ob}\,\mathrm{Cat}_\lambda$
et $\xi\in T(C)$, le foncteur
\[
\xi_{\unsure{\xi|*}}^\lambda:R_T^\lambda\longrightarrow C\tag*{(3.3)}
\]
\deleted{commute} est un $\lambda$-foncteur.

\margin{On a (pour définir (3.3))
$(3.4)\ \left\{\begin{aligned}
 R_T^\lambda\times T(C)&\longrightarrow C\\
 \deleted{\unsure{T(C)\times R_T^\lambda\to C}}&\\
 &\text{pour }C\in\mathrm{ob}\,\mathrm{Cat}_\lambda
 \end{aligned}\right.$
i.e. $(3.5)\ R_T^\lambda\longrightarrow\underline{\mathrm{Hom}}(T(C),C)$
\deleted{Notons le foncteur \unsure{Cov} correspondant \unsure{à} \unsure{commutant} aux}
qui est un $\lambda$-foncteur.}

Bien entendu, la catégorie $R_T^\lambda$ dépend de $T$ de
façon contravariante.

\margin{Quand $I$ est donné avec $T$, on a
$(\unsure{3.6|3.5})\ I\longrightarrow\mathrm{Ob}\,R_T^\lambda$ (fonctoriel en $T\in\mathrm{Ob}(\lambda\text{-}I\text{-types})$),
qui avec (3.4) redonne (2.2), et dont la donnée \deleted{\unsure{avec}} équivaut à (2.3)
i.e. (2.2) pour tt $C$.}
% margin note written sideways
\margin{Notons aussi, de façon équivalente à (3.5)
$(3.7)\ T(C)\longrightarrow\underline{\mathrm{Hom}}(R_\lambda^T,C)$
\uline{définissant}, pour $T$ variable
$(3.8)\ T_\lambda\longrightarrow\widehat{R_\lambda^T}{}^{\lambda}=(C\longmapsto\underline{\mathrm{Hom}}(R_\lambda^T,C))$
\fbox{5} \deleted{\unsure{pour l'instant}} \unsure{nonobstant} le fait que $R_\lambda^T$ soit une $\mathcal U$-catégorie \unsure{ou}}
% slanted
\margin{\deleted{\unsure{T}} $T_\lambda$ arrive (universellement ?) dans la cat. 2-cofibrée sur $(\mathrm{Cat}_\lambda)$ \uline{représentée} par $R_T^\lambda$}

\fbox{4} \uline{Théories $\lambda$-représentables}

On dit que $T$ est $\lambda$-représentable si \deleted{la catégorie} \inserted{la catégorie}
2-cofibrée $T_\lambda$ sur $(\mathrm{Cat}_\lambda)$ est 2-repr, l'objet
$R_T$ qui la représente est appelé \uline{catégorie $\lambda$-modulaire} de $T$ : \deleted{\unsure{Elle}}
\[
T(C)\xleftarrow{\simeq}\underline{\mathrm{Hom}}_\lambda(R_T,C)\tag*{(4.1)}
\]
via une \deleted{objet} $T$-structure $\lambda$-universelle
\[
\xi_T^\lambda\in T(R_T)\tag*{(4.2)}
\]
% archive sheet number: 98
\scanpage{100}
\mspage{\textcircled{3}}
On sait que $R_T$ est défini à équivalence \inserted{près} (unique à isom.
unique près), et que pour tte catégorie cofibrée $\mathcal F$ sur
$(\mathrm{Cat})_{\lambda}$, on a une équivalence de catégories (unique à isom unique près)
\[
\underline{\mathrm{Hom}}_{\mathrm{fib}/(\mathrm{Cat}_\lambda)}(T_\lambda,\mathcal F)\simeq\mathcal F(R_T)\tag*{(4.3)}
\]
Faisant $\mathcal F=T'_\lambda$, on trouve
\[
\underline{\mathrm{Hom}}_{\deleted{\unsure{\mathrm{fib}\,\mathrm{Cat}}}\lambda}(T,T')\simeq T'(R_T)\quad\text{(unique à isom unique près)}\tag*{(4.4)}
\]
et faisant en particulier $T'=\tau$, on trouve
\[
\underbrace{\mathrm{Hom}_\lambda(T,\tau)}_{\text{défn }\ R_T^\lambda}\simeq
\underbrace{\tau(R_T)}_{R_T\ \text{défn}}\qquad\text{d'où}\tag*{\deleted{(4.5)}}
\]
\[
R_T=R_T^\lambda\quad\text{(catégorie des $\lambda$-Constructions)}\tag*{(4.5)}
\]
\inserted{On \unsure{s'écrira|écrira} donc indifféremment $R_T$ ou $R_T^\lambda$. Alors (4.1)
s'interprète (à équivalence près) comme (3.8).}
\deleted{\illegible{1} Faisant \unsure{maintenant} \unsure{dans}} Dans le cas envisagé ici, la
\deleted{\illegible{1} \unsure{d'ailleurs} \unsure{soit la}} cat. des $\lambda$-constructions est une $\mathcal U$-cat.

Faisant $C=(\mathrm{Ens})$ dans (4.1), on trouve
\[
T(\mathrm{Ens})\simeq\underline{\mathrm{Hom}}_\lambda(R_T^\lambda,(\mathrm{Ens}))
\deleted{\unsure{\simeq|\cong}(R_T^\lambda)^{\unsure{\circ|o}}}\tag*{(4.6)}
\]
[si $T$ n'est pas $\lambda$-représentable, on a seulement

\fbox{5} Il faut préciser quelles sont \inserted{(mod équivalence)} les catégories
\unsure{1|2}-cofibrées sur $\mathrm{Cat}_\lambda$ (resp. munies d'un hom dans $\unsure{\lambda|\tau}_\lambda^I$)
qui proviennent des \deleted{``théories de} ``types de structures \inserted{relatifs à $\lambda$}'' -- donc
\deleted{préciser} \inserted{expliciter} une définition de la cat. ($\lambda$-types) (1.5)
\inserted{On va d'abord \deleted{\unsure{poser}} un $I$ et regarder les $\lambda$-types sur $I$.}
On va préciser \deleted{\unsure{4|les}} ``règles de formation''

a) \inserted{(Structure vide)} On a, \deleted{de $\unsure{\tau^I}$} pour $I\in\mathcal U$,
$\tau^I\in\mathrm{Ob}\,(\lambda\text{-}\inserted{I\text{-}}\mathrm{type})$

b) (\uline{\unsure{Adjonction|Rajouter} des axiomes}) \inserted{(d'isomorphie)} Si $T$ est un ($\lambda$-type) \inserted{sur $I$}, et
on se donne une \inserted{petite} partie \inserted{$A$} de $\mathrm{Fl}(S_{T,I}^{\lambda})\ (\subset\mathrm{Fl}(R_T^\lambda))$, alors
% slanted
\margin{On introduit $S_{T,I}^{\lambda}\subset R_T^\lambda$ $\hookrightarrow$ sous catégorie pleine engendrée par $I$ \deleted{\unsure{par}} \unsure{lim|lieus} de type $\lambda$ si}
la sous-catégorie 2-cofibrée pleine de $\unsure{T_A|T|T'_A}$ \deleted{formée \unsure{des}} \inserted{telle} que
$T_A(C)$ $(\deleted{\unsure{\xi|\subset}}C\in\mathrm{Ob}\,\mathrm{Cat}_\lambda)$ soit formé des $\xi\in T(C)$ tels
que (3.3) transforme $A$ en des isom, est un $\lambda$-\inserted{$I$-}type

\margin{NB peut-être faut il aussi se donner la possibilité d'égaliser des couples de flèches....}
% archive sheet number: 99
\scanpage{101}
\mspage{\textcircled{4}}
c) (Rajouter des données-flèches) Si $T$ est un $\lambda$-$I$ type \inserted{sur $I$}
et si on se donne \deleted{\unsure{des}} sur un \deleted{partie de}\inserted{petit} ensemble
$B$ et une application $B\xrightarrow{\alpha}\mathrm{Ob}\,S_{T,I}^{\lambda}\times\mathrm{Ob}\,S_{T,I}^{\lambda}$ \deleted{\unsure{et une}
\unsure{une flè}} [\unsure{i.e.|donc} une famille $B_{\rho,\sigma}$ d'ens disjoints
\unsure{du niveau|de niveaux} petits, indexée par les couples d'objets
de $S_{T,I}^\lambda$], alors la catégorie cofibrée sur $(\mathrm{Cat}_{\lambda})$
suivante est un $\lambda$-\inserted{$I$-}type :
{\small
\[
\begin{aligned}
 T_{B,\alpha}(C)={}&\deleted{\unsure{T(C)\times\prod_{\rho,\sigma}\mathrm{Hom}_{C}(\rho,\sigma)|\prod_{\rho,\sigma}\mathrm{Hom}_{C}}}\
 \text{\deleted{\unsure{catégorie}} \inserted{catégorie} des}\\
 &\left\{\begin{aligned}
 &\text{couples }\bigl(\xi,(u_{\rho,\sigma})_{\rho,\sigma\in\unsure{\{}\mathrm{Ob}\,S_{T,I}^\lambda}\bigr)\\
 &\text{avec }\xi\in\mathrm{Ob}\,T(C)\text{ et }u_{\rho,\sigma}\in\mathrm{Hom}_{C}(\rho(\xi),\sigma(\xi))\\
 &(\forall\rho,\sigma\in\mathrm{Ob}\,S_{T,I}^\lambda)
 \end{aligned}\right.
\end{aligned}
\]
}

d) \deleted{\unsure{Somme}} (Composition de structures) Si on a une
famille $(T_j)_{j\in J}$ de $\lambda$-types sur $I$, $J$ petit,
\deleted{et si on se} alors la catégorie cofibrée sur $(\mathrm{Cat}_{\lambda})$
suivante est un $\lambda$-\inserted{$I$-}type
\[
\begin{aligned}
\deleted{\prod}\ \pi_I\bigl((T_j)_{j\in J}\bigr)(C)&=\inserted{2\text{-}}\text{produit fibré des}\\
&T_j(C)\text{ sur }C^I
\end{aligned}
\]

\uline{Définition} 5.1. La 2-catégorie des $\lambda$-\inserted{$I$-}types est, par définition,
la sous-2-catégorie 2-pleine des \deleted{Cat} $\underline{2\text{-}\mathrm{Cofib}}(\deleted{\unsure{2}}\mathrm{Cat}_\lambda)/\tau^I$
\deleted{formée des \unsure{Cat. d'|cat. déf.}} \inserted{la plus petite} qui soit stable par
les op a) b) c) d)

\uline{NB} On peut en donner une construction transfinie.

\uline{Question} \unsure{5.3|5.2|5.1} \deleted{La 2-catégor} les \inserted{$I$-}cat. 2-cofibrées sur $(\mathrm{Cat})_\lambda$ qui
sont des $\lambda$-$I$-types sont elles \uline{exactement} celles qui
sont 2-représentables par un $R_{\deleted{\unsure{T}}}^{\deleted{\unsure{\lambda}}}\in\mathrm{Ob}\,\mathrm{Cat}_\lambda$, munies
de $I\longrightarrow\mathrm{Ob}\,R_{\deleted{\unsure{T}}}^{\deleted{\unsure{\lambda}}}$ telle que \deleted{systèmes de foncteurs
\unsure{sur|Ob} $R_T$} $\mathrm{Ob}\,R$
soit la plus petite partie de $\mathrm{Ob}\,R$ stable pour les $\lambda$-limites.
% archive sheet number: 100
\scanpage{102}
\mspage{\textcircled{\fbox{5}}}
5.4. L'assertion $\subset$)
\deleted{Cela} se décompose en les questions

a') $\tau^I$ est-il représentable par une catégorie $R_I$ (la
$\lambda$-catégorie libre engendrée par $I$ objets) qui soit \inserted{$\lambda$-}engendrée
(au sens de (5.3)) par $\mathrm{Im}\,(I\longrightarrow\mathrm{Ob}\,R_I)$ ?

b') Soit \deleted{\unsure{et}} $R\in\mathrm{Ob}\,\deleted{\unsure{R}}(\mathrm{Cat}_\lambda)$, $A\in\mathrm{Fl}\,R$. Y a-t-il une
$\lambda$-catégorie $R_A$ \inserted{$\in\mathrm{Ob}\,\mathrm{Cat}_\lambda$} ``universelle'' pour rendre inversibles les
flèches $\in A$, et $R_A$ est il $\lambda$-engendré par $\mathrm{Im}\,R$
% slanted
\margin{cf note pour b)}

\unsure{c'|b'|b''}$)$ Soit $R\in\mathrm{Ob}(\mathrm{Cat}_{\lambda})$, $B$ un \deleted{\unsure{ens}} petit ensemble, et
$\alpha:B\longrightarrow\deleted{\unsure{\mathrm{Ob}\,S_{T,I}\times R}}\mathrm{Ob}\,R\times\mathrm{Ob}\,R$. \deleted{\unsure{E}} Y a-t-il une
$\lambda$-cat. $R'=R_{B,\alpha}\in\mathrm{Ob}\,\mathrm{Cat}_{\lambda}$ universelle \deleted{qui} munie
d'un $\lambda$-foncteur $R\xrightarrow{i}R'$ et \unsure{d'applications|d'application}
$B\longrightarrow\deleted{\unsure{\mathrm{Ob}\,R}}\mathrm{Fl}\,R'$ telle que le composé avec $\mathrm{Fl}\,R'\to\mathrm{Ob}\,R'\times\mathrm{Ob}\,R'$
soit le composé de $\alpha$ avec $\mathrm{Ob}\,i\times\mathrm{Ob}\,i$. De plus, $R'$ est il
$\lambda$-engendré par $\mathrm{Im}\,i$ ?

d') \unsure{Soit|Soient} $(R_j)_{j\in J}$ une petite famille d'objets
de $(\mathrm{Cat}_\lambda)$, munis d'applications $I\xrightarrow{\varphi_j}\mathrm{Ob}\,R_j$. Y-a-il
une solution au pb universel d'envoyer les $R_j\xrightarrow{\psi_j}R$
$R\in\mathrm{Ob}\,\mathrm{Cat}_\lambda$, avec \deleted{des applications} \inserted{syst. transitif d'} isomorphismes entre
les composés \deleted{$\unsure{\varphi_j\circ\varphi_j}$} $\psi_j\circ\varphi_j$ ($I$ considéré comme
cat. discrète...), et $R$ est-il $\lambda$-engendré par la
réunion des images des $\mathrm{Ob}\,\psi_j$ ?

\uline{On a envie de prouver} $5.\unsure{5|6}$. Soit $S\subset\mathrm{Ob}\,R$, $R\in\mathrm{Ob}\,C_\lambda$,
supposons que $\forall R'\in\mathrm{Ob}\,C_\lambda$, et $\forall F,G\in\underline{\mathrm{Hom}}_\lambda(R,R')$,
\[
\mathrm{Hom}(F,G)\longrightarrow\prod_{s\in S}\mathrm{Hom}\bigl(F(s),G(s)\bigr)\quad\text{est injectif.}
\]
Alors $R$ est-il $\lambda$-engendré par $S$ ? (la réciproque
étant évidemment vraie...)
% archive sheet number: 101
\scanpage{103}
\mspage{\textcircled{6}}
5.6. Partons de $R\in\mathrm{Ob}\,\mathrm{Cat}_{\lambda}$, avec $I\xrightarrow{\varphi\,\unsure{\subset}} \mathrm{Ob}\,R$ tels
que $R$ soit $\lambda$-engendrée par $I$. Considérons la cat.
2-\deleted{\unsure{co}}\inserted{\unsure{co}}fibrée $T$ sur $(\mathrm{Cat}_{\lambda})$ définie par $R$, et sa structure
supplémentaire $T\xrightarrow{\varphi}\tau^I$ déduite de $\varphi$. \unsure{Prouvons|Prouverons|Prouver}
que $(T,\varphi)$ est un $\lambda$-$I$-type ? (Il faut peut-être
\margin{$|\!|\!|$}renforcer la notion de sous-catégorie $\lambda$-engendrée, en
tenant compte des flèches, pas seulement des objets... ?)

5.7. Si on laisse tomber $I$, on \deleted{trouve que}\inserted{doit présumer}
Catégorie 2-cofibrée $T$ sur $(\mathrm{Cat}_{\lambda})$ est un $\lambda$-type, si c'est
un $\lambda$-$I$-type pour \inserted{$I$ et} $\varphi:T\longrightarrow\tau^I$ convenables. Cela
signifie donc (\unsure{en|si} ce qui précède \unsure{généraliser|généralise}) que
$T$ est 2-représentable par un $R\in\mathrm{Ob}(\mathrm{Cat}_{\lambda})$, où
$\mathrm{Ob}\,R$ admet une \deleted{petite} partie qui soit $\lambda$-génératrice
(en un sens convenable à dégager).

% slanted
\margin{\uline{Question} Si $\underset{\leftarrow}{\lambda}$ et $\underset{\rightarrow}{\lambda}$ ne comportent que des diagrammes finis, en nb fini, alors une $R\in\mathrm{Ob}\,\mathrm{Cat}_\lambda$ muni d'un $I\longrightarrow\mathrm{Ob}\,R$ \unsure{provenant} d'un $\lambda$-type, est-elle une \uline{petite} catégorie ?}

\fbox{6} \uline{$\lambda$-variable} on devrait trouver \inserted{pour $\lambda\subset\lambda'$} un 2-foncteur
\[
(\lambda\text{-types})\xrightarrow{\unsure{\gamma|\eta}_{\lambda,\lambda'}}(\lambda'\text{-types})\tag*{(6.1)}
\]
\margin{$(\gamma_{\lambda,\lambda'}(T_\lambda))(C)=T_\lambda(C)$ si $C\in\mathrm{Ob}\,\mathrm{Cat}_{\lambda'}$ $(\underset{\mathrm{Ob}\,\mathrm{Cat}_{\lambda}}{\cap})$}
qui corresponde à un 2-foncteur
\[
(\mathrm{Cat}_{\lambda})\xrightarrow{\overline\gamma_{\lambda,\lambda'}}(\mathrm{Cat}_{\lambda'})\tag*{(6.2)}
\]
``$\lambda'$-catégorie enveloppante''. Il faudrait donner des
critères moyennant lesquels (6.1) est pl. fidèle.
\[
\begin{aligned}
[\ \underline{\text{Ex}}\quad\lambda&=\varprojlim\text{ finies},\quad\lambda'\ \unsure{\subset|=}\ \varprojlim\text{ quelc.}\quad(\overline\gamma(C)=\mathrm{Pro}(C)\,?)\\
\lambda&=\varinjlim\text{ finies},\quad\lambda'=\varinjlim\text{ quelc.}\quad(\overline\gamma(C)=\mathrm{Ind}(C)\,?)]
\end{aligned}
\]
\margin{? $||$} $\gamma_{\lambda,\lambda'}$ est-il tjrs \deleted{\unsure{pl}}-fidèle ?, i.e
$\underline{\mathrm{Hom}}_\lambda(T_1,T_2)\longrightarrow\mathrm{Hom}_{\lambda'}(T_1,T_2)$
est-il tjrs une équivalence pour $T_1$ et $T_2$ deux $\lambda$-types ? Sans
doute faux, mais ce doit être pl. fidèle si $\underset{\rightarrow}{\lambda}=\varnothing$,

\fbox{7} \uline{Dualité}. \uline{Opposé} d'un ``type'' $\lambda=(\underset{\leftarrow}{\lambda},\underset{\rightarrow}{\lambda})$ : on échange
les types de diagrammes utilisés pour des $\varprojlim$ resp $\varinjlim$.
Alors pour tt $\lambda$-type $T_{\deleted{\unsure{\lambda}}}$, on définit un $\lambda^o$-type $T^o$ \unsure{t.q}
\[
T^o_{\deleted{\unsure{\lambda}}}(C)\simeq T_{\deleted{\unsure{\lambda}}}(C^o)^o\tag*{(\unsure{7.1|6.3})}
\]
% archive sheet number: 102
\scanpage{104}
\mspage{\textcircled{7}}
i.e. on a une équivalence de Cat. 2-cofibrées
\[
T_{\deleted{\unsure{\lambda}}}^{o}\simeq(T_{\deleted{\unsure{\lambda}}})_{\unsure{o|\circ}}^{\unsure{o}}\quad\text{(passage à l'opposé)}\tag*{(7.1)}
\]
On a $(7.\unsure{2|3})$ \unsure{!}

On trouve ainsi une 2-équivalence
\[
(\lambda\text{-types})\simeq(\lambda^o\text{-types})\tag*{(7.3)}
\]
(et kif kif sur $I$
\[
(\lambda\text{-}I\text{-types})\simeq(\lambda^o\text{-}I\text{-types}))\tag*{(7.4.)}
\]
compatible, via les cat. représentatives, à
\[
\begin{aligned}
(\mathrm{Cat}_\lambda)&\simeq(\mathrm{Cat}_{\lambda^o})\\
C&\longmapsto C^o
\end{aligned}\tag*{(7.6)}
\]

\fbox{8} Cas où $\underset{\rightarrow}{\lambda}=\varnothing$ (structures définies via $\varprojlim$)

\inserted{$T$ désignera une $\lambda$-structure.}
8.1. Pour tout $C\in\mathrm{Ob}\,\mathrm{Cat}_{\lambda}$, les \deleted{limites}\inserted{$\varprojlim$} de type $\lambda$ existent dans $T(C)$.
plus précisément, \deleted{\unsure{les lim}} \inserted{tout type $d$} de $\varprojlim$ qui existent dans $C$
existent aussi dans $T(C)$. [Ecrire $T(C)\simeq\underline{\mathrm{Hom}}_\lambda(R_T,C)$,
et noter que \deleted{\unsure{toutes les}} l'hyp. sur $C$ implique que $\underline{\mathrm{Hom}}(R_T,C)$
\deleted{est stable} a des $\varprojlim$ de type $d$ \inserted{se calculant argt par $\underline{\text{argument}}$}, et que \deleted{\unsure{les f}} $\underline{\mathrm{Hom}}_\lambda(R_T,C)$
est \deleted{$\underline{\mathrm{Hom}}_\lambda$} est stable par lesdites $\varprojlim$. De plus, pour tt
$\unsure{r|z|\xi}\in R_T$, \deleted{\unsure{Ev}} $r_C:T(C)\longrightarrow C$ commute aux $\varprojlim$ de type $d$

8.1.1. Plus généralement ($\lambda$ \uline{quelconque}, sans néc. $\underset{\rightarrow}{\lambda}=\varnothing$)
Si \inserted{$C\in\mathrm{Ob}\,\mathrm{Cat}\,C_\lambda$ \unsure{$d$} et si} $d$ \deleted{\unsure{soit}} est un type de diagrammes tel que

a) $C$ admet des $\varprojlim$ de type $d$ (p.ex. $d\in\underset{\leftarrow}{\lambda}$)

b) dans $C$, les $\varprojlim$ de type $d$ commutent aux
$\varinjlim$ de type $\in\underset{\rightarrow}{\lambda}$

Alors \deleted{\unsure{toutes}} $T(C)$ admet des $\varprojlim$ de type $d$, et
\deleted{\unsure{si} \illegible{1} $R_T^\lambda$ alors \unsure{se}} écrivant $T(C)\simeq\underline{\mathrm{Hom}}_\lambda(R_T^\lambda,C)$,
ces limites se calculent ``argument par argument''.

% slanted
\margin{de plus si $C'$ est comme $C$ et si $C\longrightarrow C'$ commute \deleted{$T(C)\to T(C')$} aux $\varprojlim$ de type $d$, alors $T(C)\to T(C')$ \deleted{\unsure{est}}. Kif-Kif pour $T\to T'$, \unsure{alors} $T(C)\to T'(C)$ commute aux $\varprojlim$ de type \unsure{$d$}.}

toute \deleted{\unsure{flè}} $\lambda$-construction $\rho$ rel. à $T$, le foncteur
correspondant
\[
\rho:T(C)\longrightarrow C
\]
% archive sheet number: 103
\scanpage{105}
\mspage{8}
commute aux $\varprojlim$ de type $d$. Bien sûr, on a un
énoncé dual pour les $\varinjlim$

\uline{8.2.} Notons que pour \deleted{\unsure{T}} $\lambda$ quelconque \inserted{et un $\lambda$-type $T$}, il est vrai que pour
tt $\lambda$-foncteur \inserted{dans $\mathrm{Cat}_\lambda$}
\margin{(\unsure{remonte|remontez|remarque})}
\[
\deleted{\unsure{i:C\to C'}}\quad i:\deleted{\unsure{\xi}}C\longrightarrow C'\qquad\text{pl. fid. \inserted{(fid.)}}\tag*{(8.2.1)}
\]
le foncteur
\[
T(i):T(C)\longrightarrow T(C')\tag*{(8.2.2)}
\]
est pl fidèle \inserted{(fid.)}. Si $T$ est muni de $\unsure{b|\beta}:T\longrightarrow\tau^I$ \inserted{(``base'')}
\uline{fidèle}, i.e. $R=R_T^\lambda$ qui 2-représente $T$ est muni de
$\varphi:I\longrightarrow\mathrm{Ob}\,R_{\deleted{\unsure{T}}}$, alors \deleted{le carré} dans le cas pl. fidèle
\diagram{\begin{tikzcd}[ampersand replacement=\&]
 {(8.2.3)}\& T_{\lambda}(C)\arrow[r,hook]\arrow[d]\&T(C')\arrow[d]\\
 \& C^I\arrow[r,hook]\&C'{}^I
\end{tikzcd}}
est 2-cartésien [Cela signifie aussi, en termes de $R$,
que la s-catégorie pleine \deleted{de} $\underline{\mathrm{Hom}}_\lambda(R,C)$ de $\underline{\mathrm{Hom}}_\lambda(R,C')$
\deleted{s'identifie aux}\inserted{est formée des} $\lambda$-foncteurs $R\xrightarrow{\xi}C'$ tels que $\xi(\varphi(i))\in C$
pour tt $i\in I$ : en effet, comme $I$ \inserted{$\lambda$-}engendre $R$, il
s'ensuit que $\xi$ (qui commute aux lim de type $\lambda$)
envoie $R$ dans l'image essentielle de $C$.]
\inserted{En termes imagés, si $(E_i)_{i\in I}$ est une famille d'objets de $C$, ``il revient au $=$'' de se donner}
% slanted
\margin{sur ces derniers une structure de type $T$ \deleted{\unsure{dans}} sur \unsure{$C'$|$E_i$} \deleted{\unsure{fon}} et ceci pour les morphismes}

Appliquant ceci, pour $C\in\mathrm{Ob}\,\mathrm{Cat}_\lambda$, au plongt
canonique
\[
i:C\lhook\joinrel\longrightarrow\widehat C\qquad\text{qui est un $\lambda$-foncteur si }\uline{$\lambda=(\underset{\leftarrow}{\lambda},\varnothing)$}\ ;\tag*{(8.2.4)}
\]
on trouve \deleted{\unsure{que}} le foncteur
\diagram{\begin{tikzcd}[ampersand replacement=\&]
 {(8.2.5)}\& T(C)\arrow[r,dash]\arrow[d]\&T(\widehat C)\arrow[d]\\
 \&C^I\arrow[r,dash]\&\widehat C{}^I
\end{tikzcd}}
(2-cartésien)
\margin{\deleted{sup.} On suppose $\underset{\rightarrow}{\lambda}=\varnothing$}

% slanted
\margin{\uline{NB} L'hypothèse $\widehat C$ une $\mathcal U$-catégorie n'est pas essentielle en fait, car (8.2.5) ou la représentabilité de $T_\lambda$ donne à $T_\lambda(C)$ pour tte $\lambda$-catégorie $C$ (pas néc. $\in\mathrm{Ob}\,\mathrm{Cat}_\lambda$) un sens à l'équiv. près, (....) si \deleted{\unsure{les}} $T(\widehat C)$ est connu. On va essayer de le ``calculer''.}
% slanted
\margin{On suppose désormais $\lambda$ quelc.}

\deleted{On a donc pour \unsure{cela|ici},} \inserted{On suppose ici} $\widehat C$ \deleted{est une $\lambda$-catégorie} \inserted{$\in\mathrm{Ob}\,\mathrm{Cat}_{\lambda}$} \deleted{(ceci suppose que c'est} \inserted{\deleted{\unsure{soit}} que $\mathrm{ss}\,\widehat C$ est} une $\mathcal U$-catégorie \inserted{(vrai p.ex. si $C$ est $\mathcal U$-petite)} \deleted{\unsure{est de même} \unsure{comme} une \unsure{cat|catégorie}}. Donc $T(C)$ est
\deleted{On a pour tt $X\in\mathrm{Ob}$} Considérons le foncteur
contravariant canonique de $\widehat C$ dans $\mathrm{Hom}(C,(\mathrm{Ens}))$, se
restreignant à $C$ \deleted{\unsure{au sens}} moyennant $i$ est un foncteur
canonique $C^o\longrightarrow\underline{\mathrm{Hom}}_\lambda(\widehat C,(\mathrm{Ens}))$, et qui s'en va en
fait dans la \deleted{$\unsure{C\longrightarrow|\widehat C}$} ss-catégorie pleine des foncteurs
qui commutent à tt type \inserted{``petit''} de lim (ind ou proj),
donc en particulier
\[
 C^o\longrightarrow\underline{\mathrm{Hom}}_\lambda(\widehat C,(\mathrm{Ens}))\qquad X\longmapsto(F\longmapsto F(X))\tag*{(8.2.6)}
\]
Or on a
\[
\underline{\mathrm{Hom}}_\lambda(\widehat C,(\mathrm{Ens}))\longrightarrow
\underline{\mathrm{Hom}}_{\unsure{\lambda|+}}(T_{\deleted{\unsure{\lambda}}}(\widehat C),T_{\deleted{\unsure{\lambda}}}(\mathrm{Ens}))
\]
% archive sheet number: 104
\scanpage{106}
\mspage{9}
d'où \deleted{en} en composant avec (8.2.6) un foncteur canonique
\[
 C^o\longrightarrow\underline{\mathrm{Hom}}\,(T(\widehat C),T(\mathrm{Ens}))\tag*{(8.2.7)}
\]
ou, de façon équivalente, un foncteur canonique
\[
 T(\widehat C)\xrightarrow{\alpha\,\simeq}\underline{\mathrm{Hom}}\,(C^o,S_T)\deleted{\unsure{(\mathrm{Ens})}}\quad\text{où }S_T=T(\mathrm{Ens})\tag*{(8.2.8)}
\]
Si on a un ensemble d'indices de base pour $T$, on aura un
carré commutatif
\diagram{\begin{tikzcd}[ampersand replacement=\&]
 {(8.2.9)}\& T(\widehat C)\arrow[r,"{\alpha\,\simeq}"]\arrow[d,"b_{\widehat C}"']\& \underline{\mathrm{Hom}}(C^o,S_T)\arrow[d,"{\underline{\mathrm{Hom}}(C^o,b_{(\mathrm{Ens})})}"]\\
 \&\widehat C{}^I\arrow[r,dash,"\simeq" description]\&\underline{\mathrm{Hom}}(C^o,(\mathrm{Ens}))^I
\end{tikzcd}}
\margin{$\begin{aligned}
\underline{\mathrm{Hom}}_\lambda(R,\widehat C)&\simeq\underline{\mathrm{Hom}}_{\lambda\text{ en }R}(R\times C^o,(\mathrm{Ens}))\\
&\simeq\underline{\mathrm{Hom}}_{\deleted{\unsure{\lambda}}}(\unsure{C^o|C},\underline{\mathrm{Hom}}_\lambda(R,(\mathrm{Ens})))
\end{aligned}$}

\deleted{\unsure{Lorsque}}\inserted{le foncteur} $\alpha$ (8.2.8) est \inserted{tjrs} une équivalence, \deleted{alors}\inserted{donc si $\underset{\rightarrow}{\lambda}=\varnothing$} le carré 2-cartésien
\deleted{\unsure{est-il}} équivalent à
\diagram{\begin{tikzcd}[ampersand replacement=\&]
 {(8.2.10)}\&T(C)\arrow[r,hook]\arrow[d,"b_C"]\&\underline{\mathrm{Hom}}(C^o,S_T)\arrow[d,"{\varphi\mapsto(b_{i\,\mathrm{Ens}}\circ\varphi)_{i\in I}}"]\\
 \&C^I\arrow[r,hook]\&\widehat C{}^I
\end{tikzcd}}
\margin{(2 cartésien \deleted{\illegible{3}} $(\underset{\rightarrow}{\lambda}=\varnothing)$ \deleted{\unsure{pour} \illegible{1}})}

\margin{$\underset{\rightarrow}{\lambda}=\varnothing$}
Ou, en termes imagés : la donnée, pour $(E_i)\in\mathrm{Ob}\,C^I$, d'une
structure d'espèce $T$ \inserted{dans $C$}, revient à la donnée \deleted{pour chaque}
kif-kif dans $\widehat C$ i.e. à la donnée pour tt $Z\in\mathrm{Ob}\,C$ d'une
structure d'espèce $T$ sur les \deleted{ensembles}\inserted{familles}
$(E_i(Z))_{i\inZ}=(\mathrm{Hom}(Z,E_i))_{i\in I}$
variant fonctoriellement avec $Z$.

-- Lorsque $\alpha$ est une \deleted{\unsure{isom}}\inserted{équivalence}, alors on voit donc que la cat
\inserted{(pour tt $C$ tel que $\widehat C$ une $\mathcal U$-catégorie)}
\deleted{structures} $T(C)$ est connue, à équivalence près (....), lorsque
l'on connaît
\[
\left\{\begin{aligned}
&S=S_T=T(\mathrm{Ens}),\quad\text{avec les foncteurs}\\
&\beta_i:S_T\longrightarrow(\mathrm{Ens})\qquad\beta=b_{i\,\mathrm{Ens}}
\end{aligned}\right.\tag*{(8.2.11)}
\]
comme la sous-catégorie pleine de $\underline{\mathrm{Hom}}(C^o,S_T)$ formée des
$\varphi:C^o\longrightarrow S_T$ tels que les $\beta_i\,S_T$ soient représentables.

8.2.12. Si on veut une car. intrinsèque, on peut dire aussi que
c'est la sous-catégorie pleine des $\varphi:C^o\longrightarrow S_T$ tels que
$\beta_{\mathrm{Ens}}\circ\varphi$ représentable pour tt $\beta\in R_T^\lambda$. Notons que le foncteur
\[
 \beta\longmapsto\beta_{\mathrm{Ens}}\ :\ R=R_T^\lambda\longrightarrow
 \underline{\mathrm{Hom}}(\underset{S_T}{\unsure{S|S_T}},(\mathrm{Ens}))=\widehat{S^o}\quad\text{cf. (3.5)}\tag*{(8.2.13)}
\]
est aussi le foncteur évident
\[
 R\longrightarrow\underline{\mathrm{Hom}}\,(\underline{\mathrm{Hom}}_\lambda(R,\mathrm{Ens}),(\mathrm{Ens}))\tag*{(8.2.14)}
\]
Or, si on considère $S=\underline{\mathrm{Hom}}_\lambda(R,(\mathrm{Ens}))$ comme ss-catégorie
pleine de $\widehat{R^o}$, \uline{celle-ci contient $R^o$} (\inserted{identifié à} la cat. des foncteurs
représentables de $R$ dans $(\mathrm{Ens})$) donc on a \deleted{$\unsure{R^o\subset S\subset\widehat{R^o}}$}
% archive sheet number: 105
\scanpage{107}
\mspage{(10)}
\[
 R^o\subset S\subset\widehat{R^o}\tag*{(8.2.15)}
\]
\deleted{\unsure{donc}} et le foncteur
\[
 \unsure{\pi|\pi_{\xi}}:R\times S\longrightarrow(\mathrm{Ens})
\]
est donné par
\[
 \pi(c,\xi)=\deleted{\unsure{\mathrm{Hom}_{S}(c}}\ \xi(c)=\mathrm{Hom}_{\widehat{R^o}}(c,\xi)=\mathrm{Hom}_{S}(c,\xi)\tag*{(8.2.16)}
\]
i.e. \deleted{\unsure{le composé}}\inserted{l'opposé} du foncteur (8.2.13) est \deleted{induit par le}
\deleted{foncteur} composé des deux foncteurs pleinement fidèles
\[
 R^{\deleted{\unsure{o}}}\underset{(8.2.15)}{\lhook\joinrel\longrightarrow}S^o\lhook\joinrel\longrightarrow\widehat{S^o}=\underline{\mathrm{Hom}}\,(S,(\mathrm{Ens}))\tag*{(8.2.17)}
\]
donc est pl. fidèle. Ainsi $R$ est connu \inserted{à équiv. près} quand
on connaît son image dans $\widehat{S^o}$. \deleted{\unsure{En} \unsure{termes}}.

Pour résumer, si $T$ est un $\lambda$-type $(\underset{\rightarrow}{\lambda}=\varnothing)$, alors
\deleted{pour \unsure{toute} $C\in\mathrm{Ob}\,\mathrm{Cat}_{\lambda}$,} $S=T(\mathrm{Ens})$ est muni d'une sous-catégorie
pleine $\Sigma\subset S$, formée des $\varphi\in\mathrm{Ob}\,S$ tels que pour
tout $C\in\mathrm{Ob}\,\mathrm{Cat}_{\lambda}\deleted{\unsure{C}}$ et tout \deleted{\unsure{fon}} $\xi\in C_\lambda$, \deleted{\unsure{on a un fon}}
\deleted{\unsure{s'expr}} définissant un \deleted{\unsure{hom}} foncteur
\[
 \widetilde\xi:C^o\longrightarrow S\qquad X\longmapsto\mathrm{Hom}(X,\xi)
\]
le foncteur composé avec $\mathrm{Hom}_S(\varphi,-)$ soit
représentable\unsure{s} :
\[
 X\longrightarrow\mathrm{Hom}_S(\varphi,\deleted{\unsure{\widetilde\xi(X)}}\mathrm{Hom}(X,\xi))
\]
$\Sigma$ est stable dans $S$ par $\varinjlim$ \inserted{(\uline{inductives})} de type $\in\lambda$ \deleted{\unsure{\xi}} [car $C$
est stable dans $\widehat C$ par $\varprojlim$ de type $\lambda$, $\forall C\in\mathrm{Ob}\,\mathrm{Cat}_{\lambda}$]
et pour tte $C$ comme dessus, on a \unsure{deux|des} foncteurs
\[
 \Sigma^o\xrightarrow{\inserted{\lambda\text{-foncteur}}}\underline{\mathrm{Hom}}\,(T(C),C),\qquad T(C)\longrightarrow\underline{\mathrm{Hom}}_\lambda(\Sigma^o,C)
\]
On voit alors que l'on trouve ainsi une
% archive sheet number: 106
\scanpage{109}
Soit $C$ une catégorie ess petite avec $\varprojlim$ finies
$T$ l'espèce de \unsure{structure|structures} que $C$ définit
$\widehat C\simeq B_T=B$ le topos classifiant, $O_B$ la structure universelle
dans $B$ (s'identifiant à l'inclusion $C\longrightarrow\widehat C$)
$E$ un topos

\diagram{\begin{tikzcd}[ampersand replacement=\&,column sep=2pt]
 O_E\arrow[d,Leftright]\&{}\in\mathrm{Ob}\,T(E)\simeq\mathrm{Ob}\,\underline{\mathrm{Hom}}_{\unsure{\mathrm{sex}|\mathrm{lex}}}(C,E)\ni f_0\\
 \deleted{\unsure{f}}\&{}=f:\unsure{E'|E}\longrightarrow B_T\quad\text{le morphisme correspondant}
\end{tikzcd}}
\deleted{$\unsure{\xi|f_0}$ la structure} donc
\[
 O_E\simeq f_T^*(O_B)\qquad f_0=f^*|C
\]

$U\in\mathrm{Ob}\,E$, $X\in\mathrm{Ob}\,B$ \inserted{(p. ex. $X\in\mathrm{Ob}\,C$)}, on va définir une
\deleted{bijection}\inserted{application}
\[
 \mathrm{Hom}_E(U,f^*X)\xrightarrow{\deleted{\unsure{\sim}}}
 \mathrm{Hom}_{T(\mathrm{Ens})\inserted{=\mathrm{Ind}(C^o)}}\bigl(O_B(X),O_E(U)\bigr)\tag*{(1)}
\]
\deleted{\unsure{lim}}
Où $O_E(U)\in T(\mathrm{Ens})$ est l'image de $O_E$ \inserted{$\in T(E)$} par le
foncteur $Z\longmapsto\mathrm{Hom}_E(U,Z):E\longrightarrow(\mathrm{Ens})$ (qui
commute aux $\varprojlim$) et où on définit de $=$
$O_B(X)$ dans $\widehat C$. En effet, on a
{\scriptsize
\[
 \mathrm{Hom}_E(\unsure{M|U},f^*(X))\longrightarrow
 \mathrm{Hom}_{T(\mathrm{Ens})}\bigl(O_E(f^*(X)),O_E(U)\bigr)
 \longrightarrow\mathrm{Hom}_{T(\mathrm{Ens})}\bigl(\underline{O}_B(X),O_E(\unsure{M|U})\bigr)
\]
}
car on a un morphisme de $T(\mathrm{Ens})$
\[
 O_B(X)\longrightarrow O_E(f^*(X))
\]
\diagram{\begin{tikzcd}[ampersand replacement=\&,row sep=small]
 C\arrow[r]\arrow[dr,"f_0"']\&\widehat C\arrow[r,"h_X",""{name=upper,below}]\&(\mathrm{Ens})\\
 \&E\arrow[ur,"h_{f^*(X)}"',""{name=lower,above}]\arrow[from=upper,to=lower]\&{}
\end{tikzcd}
\quad
\begin{tikzcd}[ampersand replacement=\&,row sep=small]
 Y\arrow[r,mapsto]\&\mathrm{Hom}_{\widehat C}(X,Y)\arrow[dr]\&{}\\
 Y\arrow[rr,mapsto]\&{}\&\mathrm{Hom}_E(f^*(X),f^*(Y))
\end{tikzcd}}
% archive sheet number: 108
\scanpage{110}
Prouvons que si $X\in\mathrm{Ob}\,C$, l'application précédente est
une bijection.

\inserted{D'ailleurs, le $2^{\mathrm e}$ membre s'exprime comme
$\mathrm{Hom}_{\mathrm{Ind}(C^o)}(X^o,O_E(U))$ pour $U$ fixé; la catégorie des couples $(X^o;\xi)$ $(\xi\in\mathrm{Hom}(X^o,O_E(U)))$
est la catégorie d'indices naturelle définissant le Ind-objet $O_E(U)$ de $\mathrm{Ind}(C^o)$....}
\[
\begin{aligned}
 O_B(X):{}&C\longrightarrow(\mathrm{Ens}) &Y&\longmapsto\mathrm{Hom}_C(X,Y)=\mathrm{Hom}_{C^o}(Y,X)\\
 O_E(U):{}&&Y&\longmapsto\mathrm{Hom}_E(\unsure{U|M},f_0(Y))
\end{aligned}
\]
\[
 \mathrm{Hom}(\underline{O}_B(X),O_E(U))=O_E(U)(X)=\mathrm{Hom}_E(U,f_0(X))
\]

\centerline{\rule{5em}{0.4pt}}

On va en profiter pour calculer $f^*(\mathcal F)$, où
$\mathcal F$ est un champ sur $B$. Il revient au $=$ de
se donner la restriction de $\mathcal F$ au \deleted{\unsure{C}} site de
définition $C$, i.e. de donner une catégorie fibrée sur $C$,

[\uline{NB} si on a un sous-topos \inserted{$B'$} de $B=\widehat C$, correspondant à une
topologie $\tau$ sur $C$, et à une sous-espèce de structure
$T'$ de $T$, \deleted{alors} si on suppose que $O_E$ est de
cette espèce i.e. que $f$ se factorise \inserted{en $f'$} par $B'$, i.e.
que $f_0$ est continu pour $\tau$, alors la donnée d'un
champ \inserted{$\mathcal F'$} sur $\underline B'$ revient de $=$ à la donnée d'un
champ $\mathcal F'_0$ sur $B$ muni de $\tau$, et \unsure{]} notre calcul
va donner $f'{}^*(\mathcal F')$ en termes de $f_0$ et $\mathcal F'$).

On \deleted{trouve que}\inserted{commence} par \deleted{que la catégorie fibrée}\inserted{noter que la catégorie fibrée}
\deleted{sur $E$} $\mathcal F'_0$ \inserted{sur $C$} se prolonge en Catégorie cofibrée sur
$\mathrm{Pro}(C)\simeq T(\mathrm{Ens})^o$, \deleted{soit} $\overline{\mathcal F'_0}$, donc $\overline{\mathcal F'_0}(O_E(U))$
est défini [mais covariant en l'objet $(\ )\in T(\mathrm{Ens})$.
Pour $U$ variable, $U\longmapsto\overline{\mathcal F'_0}(O_E(U))$ donne donc une
catégorie fibrée sur $U$ [image inverse de $\overline{\mathcal F'_0}$ par $\underline O_E:E^o\to T(\mathrm{Ens})\underset{\unsure{\simeq|=|\cong}}{\vphantom A}$
$\mathrm{Pro}(C)^o$
On prend le \deleted{catégorie}\inserted{champ} associé.
% archive sheet number: 109
\scanpage{112}
\uline{Topos modulaire pour des structures algébriques
\deleted{\unsure{définies}} sur objets \inserted{fixés} d'un topos fixé \deleted{\unsure{E}} $X$}

\deleted{\unsure{On}} On peut se donner des objets $E_i\in X$, $i\in I$
et \deleted{de s'inté} une espèce de structure $T$ avec
topos modulaire $B$, \deleted{sur objets} avec donnée
``objets de base'' correspondants : des $b_i\in B$, $i\in I$

On s'intéresse donc, pour tt topos $X'$ sur $X$,
\inserted{(ayant comme images inverses des $E_i$ les $E'_i$)}
aux structures d'espèce $T$ \deleted{dans} $X'$ \deleted{qui}
(i.e. morphismes $f:X'\longrightarrow B$) \deleted{tels que} avec
isom donnés $f^*(b_i)\simeq E'_i$. La catégorie
\diagram{\begin{tikzcd}[ampersand replacement=\&,row sep=small]
 B\arrow[d,"b"]\&X\arrow[dl,"E"]\\
 \widehat{(\mathrm{Ens\,f})^o}{}^{I}=B_I\&{}
\end{tikzcd}}
de ces objets sur un topos $X'$
fixé, pour $X'\longrightarrow X$ \uline{et} $X'\longrightarrow B$
variable, est le produit 2-fibré de $\underline{\mathrm{Hom}}_{\mathrm{top}}(X',X)$ et
$\underline{\mathrm{Hom}\,\mathrm{top}}(X',B)$ sur $\underline{\mathrm{Hom}\,\mathrm{top}}(X',B_I)$
Donc on cherche un produit
2-fibré $B\times_{B_I}X$, dans la 2-catégorie (Top)
(On n'a pas encore prouvé que cela marche)
\deleted{Cela doit donc être une sous-\unsure{espèce}}

Supposons qu'il s'agisse d'une $T$ d'une
structure définie par $\varprojlim$ finies, avec cat. modulaire
$R$ (donc $B=\widehat R$ \unsure{$\}$}, $I\lhook\joinrel\longrightarrow\mathrm{Ob}\,R$) Donc on
s'intéresse aux foncteurs exact à g.
$R\longrightarrow X'$ transformant les $r_i$ en les $E'_i$.
Soit $\unsure{\mathscr S|S|\mathcal S}_{\deleted{\unsure{X}}}$ \deleted{la catégorie}\inserted{le champ sur $X$} formée des structures
de ce type dans $(X,(E_i))$, \inserted{et les fonctrs induits} i.e. des functs ex
à g. $f:R\longrightarrow X$ avec $f(r_i)\simeq E_i$ donnés, et
% slanted
\margin{champ en \unsure{cat.} à choisir (associé à un faisc. en cat. donné ?)}
\deleted{et soit Considérons la}
soit. $\unsure{\Phi|I|\underline I}_X$ le champ \inserted{en topos} des flèches de $\deleted{\unsure{B}}X$.
% archive sheet number: 111
